\documentclass[11pt,a4paper]{article}
\usepackage{templates/modeling-notes}
\hypersetup{pdftitle={Mathematical modeling -- Lab 2: linear programming}}
\title{\textbf{Mathematical Modeling}\\[2mm]\large Lab 2: distributions and courier deliveries}
\author{Personal study notes}
\date{}
\begin{document}
\maketitle
\tableofcontents

\section{From a statement to solver matrices}
Both tasks ask us to choose quantities subject to linear restrictions. First
define the variables, derive the objective and constraints, and check feasibility.
Only then assemble the numerical arrays. We use column decision vectors and the
convention
\[
  \min f^Tv,\qquad Av\leq b,\qquad A_{\mathrm{eq}}v=b_{\mathrm{eq}}.
\]
Every row is one restriction; every column belongs to one decision variable.
Bounds can be written as matrix rows:
\[
  \ell\leq v\leq u
  \quad\Longleftrightarrow\quad
  \underbrace{\begin{bmatrix}-I\\I\end{bmatrix}}_{A}v
  \leq\underbrace{\begin{bmatrix}-\ell\\u\end{bmatrix}}_b.
\]
The lower bound changes sign because $v_i\geq\ell_i$ becomes
$-v_i\leq-\ell_i$. The MATLAB scripts below put all bounds into these rows
and pass empty lower/upper-bound arguments to \texttt{linprog}. Using its
\texttt{lb}, \texttt{ub} arguments instead describes the same feasible set.

\section{Group 1: choosing a probability distribution}\label{sec:probability}
\subsection{Task statement}
\begin{otherlanguage}{polish}
  Dana jest zmienne losowe $X$ oraz $Y$ o dziedzinie
  $[5]=\{1,2,3,4,5\}$ przyjmujące wartości zgodnie z macierzami
  przedstawionymi na dole tego pliku.
  Naszym celem jest znalezienie rozkładu prawdopodobieństwa
  $P=(p_i)_{i\in[5]}$ takiego, który minimalizuje wartość:
  \begin{enumerate}[label=\alph*)]
    \item $E_P[X]$ (przy dodatkowym założeniu, że każde prawdopodobieństwo
          $p_i$ (dla $i\in[5]$) jest co najwyżej równe $0.3$),
    \item $z$ ograniczającą z góry jednocześnie wartości: $E_P[x]$ oraz $E_P[Y]$.
  \end{enumerate}
  Zaprojektuj i zaimplementuj dwa modele programowania liniowego, które znajdą
  optymalne rozkłady prawdopodobieństwa dla powyższych podpunktów. Ile wynoszą
  wartości oczekiwane w każdym z tych przypadków? W każdym podpunkcie,
  zwizualizuj uzyskane rozkłady prawdopodobieństwa przy pomocy wykresów słupkowych
  (przydatna może być funkcja \texttt{bar}).
\end{otherlanguage}
The supplied value arrays are
\[
  \begin{array}{c|rrrrr}
    i&1&2&3&4&5\\\hline
    X(i)&1.2&2.1&-0.4&1.2&-0.5\\
    Y(i)&1.2&-0.6&2.4&-0.1&0.5
  \end{array}.
\]

\subsection{Variables and expectations}
The outcomes and their values are fixed. The decisions are the probabilities
$p=(p_1,\ldots,p_5)^T$, with $p_i\geq0$ and $\sum_i p_i=1$.
Use the value vectors
\[
  \xi=(1.2,2.1,-0.4,1.2,-0.5)^T,\qquad
  \eta=(1.2,-0.6,2.4,-0.1,0.5)^T.
\]
Then
\[
  E_P[X]=\xi^Tp,\qquad E_P[Y]=\eta^Tp.
\]
These are weighted averages, so they are linear in the unknown probabilities.
The $0.3$ cap is an extra condition in part (a) only.

\subsection{Part (a): minimize one expectation}
The complete LP is
\[
  \min\xi^Tp,\qquad
  \begin{bmatrix}-I_5\\I_5\end{bmatrix}p
  \leq\begin{bmatrix}0_5\\0.3\,\mathbf1_5\end{bmatrix},
  \qquad\mathbf1_5^Tp=1.
\]
To lower the expectation, place as much probability as allowed on the smallest
values: $X(5)=-0.5$ and $X(3)=-0.4$ each receive $0.3$.
The remaining $0.4$ goes to outcomes 1 and 4, both with value $1.2$.
Outcome 2 has a larger value and receives zero. Thus all optimal distributions are
\[
  p=(t,0,0.3,0.4-t,0.3)^T,\qquad 0.1\leq t\leq0.3.
\]
The interval ensures that both $p_1$ and $p_4$ remain between zero and $0.3$.
For a direct lower-bound certificate, eliminate $p_1+p_4$ using the sum:
\[
  E_P[X]=1.2+0.9p_2-1.6p_3-1.7p_5
  \geq1.2-1.6(0.3)-1.7(0.3)=0.21.
\]
Every displayed distribution attains this bound. Its other expectation is
\[
  E_P[Y]=0.83+1.3t\in[0.96,1.22].
\]
For example, $p=(0.1,0,0.3,0.3,0.3)^T$ gives
$E_P[X]=0.21$ and $E_P[Y]=0.96$. Different optimal bar charts can therefore
be equally correct: part (a) minimizes $E_P[X]$, not $E_P[Y]$.

\subsection{Part (b): minimize a common upper bound}
For a fixed distribution, the smallest number bounding both expectations is
$\max\{E_P[X],E_P[Y]\}$. Introduce $z$ as a sixth decision variable and solve
\[
  \min z,\qquad \xi^Tp-z\leq0,\qquad\eta^Tp-z\leq0,
  \qquad p\geq0,\quad\mathbf1_5^Tp=1.
\]
For $v=(p_1,\ldots,p_5,z)^T$, the solver arrays are
\[
  f=\begin{bmatrix}0_5\\1\end{bmatrix},\quad
  A=\begin{bmatrix}\xi^T&-1\\\eta^T&-1\\-I_5&0_5\end{bmatrix},\quad
  b=0_7,\quad A_{\mathrm{eq}}=\begin{bmatrix}\mathbf1_5^T&0\end{bmatrix},
  \quad b_{\mathrm{eq}}=1.
\]
Here $A$ has seven rows and six columns. Its last five rows restrict the
probabilities, leaving $z$ unrestricted in sign. The $0.3$ cap is absent.

\paragraph{Hand solution and a global certificate.}
Outcome 5 favors $X$; outcome 2 favors $Y$. Mixing just these two, set
$p_2=t$, $p_5=1-t$. Their expectations are
\[
  E_P[X]=-0.5+2.6t,\qquad E_P[Y]=0.5-1.1t.
\]
Balancing them gives $3.7t=1$, hence
\[
  p^*=(0,10/37,0,0,27/37)^T,\qquad
  E_{P^*}[X]=E_{P^*}[Y]=z^*=15/74\approx0.202703.
\]
Balancing a chosen pair alone would not rule out a better mixture involving
other outcomes. To prove global optimality, observe that
\[
  \frac{11\xi+26\eta}{37}
  =\left(\frac65,\frac{15}{74},\frac{58}{37},\frac{53}{185},\frac{15}{74}\right)^T
  \geq\frac{15}{74}\,\mathbf1_5.
\]
For any feasible distribution and upper bound,
\[
  z\geq\frac{11}{37}E_P[X]+\frac{26}{37}E_P[Y]
  \geq\frac{15}{74}\sum_i p_i=\frac{15}{74}.
\]
Our distribution attains this bound. Equality requires zero mass at outcomes
1, 3 and 4 and equal expectations, so this optimum is unique.

\subsection{MATLAB implementation and comparison}
Run \href{run:matlab/labs/lab_02_probability.m}{\texttt{lab\_02\_probability.m}}
from the repository root:
\begin{lstlisting}[language=Matlab]
run('matlab/labs/lab_02_probability.m')
\end{lstlisting}
The script prints both expectations and the probability vectors, draws one bar
chart for each part, and checks feasibility and objective values with absolute
tolerance $10^{-7}$. It uses solver constraint and optimality tolerances
$10^{-9}$. For part (a), it checks the optimal family rather than requiring
one particular vector. For part (b), it checks the unique vector above.

\section{Group 5: deliveries along a fixed courier route}\label{sec:courier}
\subsection{Task statement}
\begin{otherlanguage}{polish}
  W pewnej firmie kurierskiej kurier przewozi paczki. Swoją trasę rozpoczyna
  w centrum logistycznym, a następnie podąża zgodnie z trasą przedstawioną
  na poniższym rysunku (centrum logistyczne oznaczono jako 0), kończąc w punkcie
  0 po odwiedzeniu wcześniej wszystkich punktów.
\end{otherlanguage}
\begin{figure}[H]
  \centering
  \begin{tikzpicture}[every node/.style={font=\small}]
    \node[draw,circle] (zero) at (0,2) {0};
    \node[draw,circle] (one) at (3,2) {1};
    \node[draw,circle] (two) at (3,0) {2};
    \node[draw,circle] (three) at (0,0) {3};
    \draw[model flow] (zero)--node[above] {$2$} (one);
    \draw[model flow] (one)--node[right] {$3$} (two);
    \draw[model flow] (two)--node[below] {$1$} (three);
    \draw[model flow] (three)--node[left] {$1.5$} (zero);
  \end{tikzpicture}
  \caption{The prescribed route and its edge cost weights. Only delivery
    quantities are decisions; the route is fixed.}
\end{figure}
\begin{otherlanguage}{polish}
  Kurier rozpoczyna podróż z ładunkiem $M$ kilogramów. Następnie, w każdym
  kolejnym punkcie podróży zostawia pewną liczbę kilogramów, przy czym:
  \begin{itemize}
    \item w punkcie 1 musi zostawić przynajmniej 2 kilogramy, a maksymalnie 6 kilogramów,
    \item w punkcie 2 musi zostawić przynajmniej 4 kilogramy, a zostawić maksymalnie 7 kilogramów,
    \item w punkcie 3 musi zostawić przynajmniej 2.5 kilograma, a zostawić maksymalnie 5 kilogramów.
  \end{itemize}
  Zaprojektuj odpowiedni model programowania liniowego, który wyznaczy, ile
  kilogramów musi zostawić w każdym z punktów tak, aby spełnić warunki i
  zminimalizować koszt podróży, który liczony jest następująco:
  \begin{itemize}
    \item koszt przejechania krawędzi to jej waga przemnożona przez wagę wiezionego ładunku,
    \item koszt całej podróży to sumaryczny koszt przejechanych krawędzi.
  \end{itemize}
  \begin{sloppypar}
    Rozwiąż zaprojektowany model przy pomocy funkcji \texttt{linprog} kolejno
    dla wartości $M=7,8,9,10,11$. Ponadto, przy pomocy wykresu słupkowego
    zwizualizuj uzyskane rozwiązania dla każdego punktu na trasie i każdej
    wartości $M$ (przydatna może być funkcja \texttt{bar}).
  \end{sloppypar}
\end{otherlanguage}

\subsection{Variables, loads, and the objective}
Let $x_i$ be kilograms delivered at point $i$. The initial mass $M$ is a given
parameter for each run. Masses are continuous. The load decreases after each
delivery, so the objective is the sum of the following costs:
\begin{center}
  \begin{tabular}{cll}\toprule
    Edge & Load on departure (kg) & Edge cost\\\midrule
    $0\to1$ & $M$ & $2M$\\
    $1\to2$ & $M-x_1$ & $3(M-x_1)$\\
    $2\to3$ & $M-x_1-x_2$ & $M-x_1-x_2$\\
    $3\to0$ & $M-x_1-x_2-x_3$ & $1.5(M-x_1-x_2-x_3)$\\\bottomrule
  \end{tabular}
\end{center}
Thus the total trip cost is
\[
  C_M(x)=7.5M-5.5x_1-2.5x_2-1.5x_3.
\]
Delivering one kilogram at point 1 saves the weights of all remaining edges,
$3+1+1.5=5.5$. The analogous savings at points 2 and 3 are $2.5$ and $1.5$.
This explains the negative coefficients. The source does not specify a currency
unit, so costs are reported in the units induced by its edge weights.

\subsection{The complete matrix model}
With $x=(x_1,x_2,x_3)^T$, minimize $f^Tx$ using
\[
  f=\begin{bmatrix}-5.5\\-2.5\\-1.5\end{bmatrix},\qquad
  \underbrace{\begin{bmatrix}
      1&1&1\\
      -1&0&0\\0&-1&0\\0&0&-1\\
      1&0&0\\0&1&0\\0&0&1
    \end{bmatrix}}_{A\in\mathbb R^{7\times3}}x
  \leq
  \underbrace{\begin{bmatrix}M\\-2\\-4\\-2.5\\6\\7\\5\end{bmatrix}}_b.
\]
There are no equality constraints. The first row prevents delivery of more
than the initial load. Because all deliveries are nonnegative, this also
ensures that every earlier edge load is nonnegative. The next three rows
are the minimum deliveries; the last three are the maximum deliveries.
The condition $x_1+x_2+x_3\geq0$ is already implied and needs no extra row.
The statement allows returning with undelivered load; an equality
$\sum_i x_i=M$ need not be imposed as a modeling assumption.

The term $7.5M$ is constant within a run, so omitting it from the solver
objective preserves the optimizer. Add it back to the returned objective
value when reporting the cost of the trip, particularly when comparing runs
with different $M$.

\subsection{Hand solution and feasibility}
The minimum deliveries total $2+4+2.5=8.5$ kg. Thus $M=7$ and $M=8$
are infeasible. For $M=9,10,11$, start with these minimum deliveries and
allocate the remaining $M-8.5$ kg. Each extra kilogram saves most when
delivered at point 1. Its upper bound allows four extra kilograms, more
than the $0.5$, $1.5$, or $2.5$ kg available in these cases. Hence
\[
  x^*(M)=(M-6.5,4,2.5)^T,\qquad M\in\{9,10,11\}.
\]
For a global certificate, write $S=x_1+x_2+x_3$. Every feasible point satisfies
\[
  5.5x_1+2.5x_2+1.5x_3
  =5.5S-3x_2-4x_3\leq5.5M-3(4)-4(2.5)=5.5M-22.
\]
Consequently $C_M(x)\geq2M+22$, with equality at the stated delivery vector.
Equality requires $S=M$, $x_2=4$, and $x_3=2.5$, so each feasible case here
has a unique optimizer.
\begin{center}
  \begin{tabular}{crrrr}\toprule
    $M$ (kg) & $x_1$ (kg) & $x_2$ (kg) & $x_3$ (kg) & Minimum trip cost\\\midrule
    7 & \multicolumn{4}{c}{Infeasible}\\
    8 & \multicolumn{4}{c}{Infeasible}\\
    9 & 2.5 & 4 & 2.5 & 40\\
    10 & 3.5 & 4 & 2.5 & 42\\
    11 & 4.5 & 4 & 2.5 & 44\\\bottomrule
  \end{tabular}
\end{center}
All available load is delivered at these optima because each delivery saves
cost and the upper bounds do not prevent it. This is a conclusion for these
inputs, rather than an extra equality added to the task.

\subsection{MATLAB implementation and comparison}
Run \href{run:matlab/labs/lab_02_courier.m}{\texttt{lab\_02\_courier.m}}:
\begin{lstlisting}[language=Matlab]
run('matlab/labs/lab_02_courier.m')
\end{lstlisting}
The script runs all five requested loads, checks infeasibility for 7 and 8,
and compares the remaining cases with the hand solution. It independently
recomputes the full trip cost from the four edge loads. Solver tolerances are
$10^{-9}$; comparison tolerance is $10^{-7}$. The grouped bar chart retains
ticks for every $M$ and uses missing values for infeasible cases: zero bars
would incorrectly suggest a feasible zero-delivery solution.

\section{Sources and editorial remarks}
Task sources: Hubert Grochowski, \texttt{Sam2Gr1.pdf} and
\texttt{Sam2Gr5.pdf}, one page each, headed Lab 2, academic year 2026/2027.
Originals are at \path{../Labs/tasks/} relative to the repository root.
The mathematical task wording is retained in Polish above, with the supplied
arrays typeset as a table and the courier diagram redrawn. Administrative
submission instructions are omitted. The lowercase $x$ in the probability
statement's $E_P[x]$ is interpreted as the named random variable $X$.

The user's completed courier submission is
\path{../Labs/Lab2/RadoslawGlasekLab2.m}; the user reports full lab credit.
\path{../Labs/Lab2/Lab2.m} is a friend's probability solution and was consulted
as a comparison, not copied into the companions. The Lab 3 file is a starter
copy, not evidence of completed Lab 3 work. The original files are preserved.
The matrix-only formulations, hand certificates and companion scripts are
independent study additions.

Exact rational checks confirm the displayed supporting coefficients, minimax
mixture and courier costs. The agent attempted to run both companions through
the installed MATLAB launcher on 25 October 2026; it exited with code 1 and
no diagnostic before producing output. Actual MATLAB execution and plot
rendering therefore remain unverified. Expected results and acceptance
tolerances are given above; execution status is also recorded in
\path{docs/matlab-verification.md}.
\end{document}
